\bookchapter{第七章：前沿功能全解析}

\chapterintro{科技日新月異，手機之功亦不斷進化。衛星通、無網通、車禍檢、星閃之技……此等聞之甚幻之功，實已入吾等之生。令吾等同探此等前沿之技。}

\booksection{一、衛星通：無信亦能求救}

衛星通乃近年手機之域一大突。當汝於荒郊野嶺、深山老林，手機全無信時，衛星通或乃救命之草。

\booksubsection{工作之理：}

手機直與軌道上之通信衛星建連，發求救之信。今主支持發短信與位之信，不能打電話。

\booksubsection{支持之機型：}

\begin{itemize}
  \item \textbf{蘋果}：iPhone 14系列及以後，支持通過衛星發送緊急SOS。iPhone 17系列進一步支持衛星發送iMessage與短信，且擴展矣更多地區之覆蓋。
  \item \textbf{華爲}：Mate 50 系始支持北斗衛星消息（惟文），Mate 60 Pro 系進一步支持天通衛星通話（可打電話），Mate 70 系衛星通之能進一步強。
  \item \textbf{小米/OPPO/vivo}：2024年起旗艦機型續支持北斗衛星消息之功。
  \item \textbf{發展之勢}：衛星通正自高端旗艦專之功向下普及，未來中端機型亦有望支持。中國北斗衛星系統於國產手機上之支持至廣。
\end{itemize}

\booksubsection{使用之場：}

\begin{itemize}
  \item 戶外探險、登山遇險。
  \item 戶外探險、登山遇險。
  \item 海上遇險。
\end{itemize}

\booksubsection{使用之法：}

\begin{enumerate}
  \item 保於開之野，能見天。
  \item 保於開之野，能見天。
  \item 依屏之指，對衛星之方。
  \item 發預設之求救之信與位。
\end{enumerate}

\begin{tcbinfo}{小知}
小知

\end{tcbinfo}

\booksection{二、無網通：近距通訊之新方}

“無網通”指於無移動網與Wi-Fi之況下，兩臺手機之間直通訊之技。

\booksubsection{技之理：}

用藍牙或專用之無線電之頻，於短距內（通常數百米至數公里）建點對點之連。

\booksubsection{應用之場：}

\begin{itemize}
  \item \textbf{戶外之動}：徒步、露營時隊友間之通。
  \item \textbf{災之援}：地震、洪水後通中時。
  \item \textbf{大型之動}：演唱會、體育賽事人多信差時。
\end{itemize}

\booksubsection{支持之台：}

\begin{itemize}
  \item \textbf{蘋果}：iOS 17之"NameDrop"與"離線找友"之功。
  \item \textbf{華爲}：暢連App支持無網通。
  \item \textbf{小米}：小米妙享中心之部分功。
\end{itemize}

\booksection{三、車禍之檢：要時救命}

車禍之檢乃用手機之感器檢重車禍，並自動撥緊電之功。

\booksubsection{工作之理：}

手機經加速度計、陀螺儀、氣壓計、麥克風等多感器，合判是否生重撞。檢到車禍後，手機將：

\begin{enumerate}
  \item 彈倒計時之界，問汝是否無恙。
  \item 彈倒計時之界，問汝是否無恙。
  \item 向緊聯人發汝之位。
\end{enumerate}

\booksubsection{支持之機型：}

\begin{itemize}
  \item \textbf{蘋果}：iPhone 14 系及以後。
  \item \textbf{安卓}：部分高端機型續支持。
\end{itemize}

\booksubsection{注意之事：}

\begin{itemize}
  \item 功須開方能用。
  \item 功須開方能用。
  \item 或誤觸（如過山車、劇動），即取可也。
\end{itemize}

\booksection{四、星閃之技：下代之無線連}

星閃（SparkLink）乃華爲牽頭研之新代無線通訊之技，目標乃於某場替藍牙與Wi-Fi。其於關鍵之標對標甚超現之方：空口之速可達九百Mbps（遠超藍牙之五十Mbps），延遲低至二十$\mu$ s（藍牙約十五ms），單網支持四千九十六設備同連（藍牙通常八以內）。

\booksubsection{技之特：}

\begin{itemize}
  \item \textbf{超低之延}：端至端之延惟二十微秒（零點零二ms），較藍牙低兩量級，遊戲與樂之同步幾乎無可感。
  \item \textbf{高之速}：單鏈峯值之速九百Mbps，乃藍牙音頻之數十倍，可傳無損音頻甚高清視頻。
  \item \textbf{低之耗}：待機之耗與藍牙當，而傳之效更高，單位數據之能耗更低。
  \item \textbf{廣之覆加高之並}：連之距較藍牙遠五成以上，單網支持上千設備同在。
\end{itemize}

\booksubsection{應用之場：}

\begin{itemize}
  \item \textbf{無線之機}：更低之延，玩遊戲音畫同步。
  \item \textbf{智能之表}：更穩之連。
  \item \textbf{智能之家}：更多設備同連。
  \item \textbf{車載之聯}：更流之手機-車機之連。
\end{itemize}

\booksubsection{支持之設備：}

今主要爲華爲生態之設備，如華爲 Mate 60 系、部分華爲耳機與表。未來將有更多之品牌支持。

\textbf{技之望}\par

技之望

\begin{tcbpractice}{往略檢之：汝手機是否支持此等功}
\begin{enumerate}
  \item 往設置中搜"衛星"，觀有無衛星通相關之擇。
  \item 搜"車禍檢"或"緊急 SOS"，檢是否開。
  \item 若汝手機乃華爲，觀是否支持星閃設備之連。
  \item 解此等功之用法，要時或能派上用場。
\end{enumerate}

\end{tcbpractice}
